\chapter{Implementation}

\section{Development Environment}
Building the Quantum Virtual Machine required a development environment tuned for both numerical linear algebra and compiler construction. We chose Python 3.10+ as the core language because it provides structural pattern matching, improved type hints, and mature scientific computing libraries.

\textbf{Core Development Toolchain:}
\begin{itemize}[leftmargin=*]
    \item \textbf{Language Environment:} Python 3.10+ running in a dedicated virtual environment (\texttt{.venv}) to isolate dependency versions.
    \item \textbf{IDE and Tooling:} Visual Studio Code with the Python, Pylance, and Git extensions. Pylance provided real-time static type checking, catching subtle shape mismatches in NumPy tensor operations during development.
    \item \textbf{Version Control:} Git hosted on GitHub. We used feature branches for compiler components (e.g., \texttt{feature/sabre-router}, \texttt{feature/mps-simulator}), ensuring that core simulator tests never broke while testing new transpilation heuristics.
    \item \textbf{Dependency Management:} Configured via \texttt{pyproject.toml} and locked in \texttt{requirements.txt}. We kept runtime dependencies minimal: NumPy \cite{numpy2020}, Lark \cite{lark2023}, FastAPI \cite{fastapi2023}, and Matplotlib \cite{matplotlib2007}.
    \item \textbf{Automated Test Suite:} Automated through \texttt{pytest}, containing 215 unit, integration, and algorithmic verification tests.
\end{itemize}

\section{Core Module Implementation}

\subsection{Module 1: Parser Layer}
The parser layer ingests quantum program source text and constructs the system's intermediate representation. It supports three distinct formats:

\textbf{OpenQASM 3.0 Parser (\texttt{src/qvm/qasm3\_parser.py}):}
We implemented this parser using Lark \cite{lark2023}, a modern pure-Python parsing toolkit. The language grammar is defined declaratively in \texttt{src/qvm/qasm3.lark} using Extended Backus-Naur Form (EBNF). Lark generates a fast LALR(1) shift-reduce parser that produces a raw Abstract Syntax Tree (AST). 

Because OpenQASM 3.0 supports both classical and quantum variables, we process the AST in two sequential passes:
\begin{enumerate}[leftmargin=*]
    \item \textbf{Declaration Pass:} Walks the tree to record qubit registers (\texttt{qubit[n] q;}) and classical bit arrays (\texttt{bit[n] c;}). This allocates contiguous integer indices and creates name-to-index lookup maps.
    \item \textbf{Instruction Pass:} Iterates through gate calls, measurements, conditionals, and loops. Gate identifiers are mapped to internal operation types, angle parameters are evaluated, and classical branch conditions are linked to classical memory addresses.
\end{enumerate}

\textbf{OpenQASM 2.0 Parser (\texttt{src/qvm/parser.py}):}
A lightweight parser for legacy OpenQASM 2.0 files. It tokenizes standard header declarations (\texttt{OPENQASM 2.0;}, \texttt{include "qelib1.inc";}), register statements, and canonical gate applications.

\textbf{JSON Gate List Parser (\texttt{src/qvm/parser.py}):}
A direct dictionary deserializer for programmatically generated circuits. It accepts JSON payloads containing explicit arrays of gate names, target qubit indices, and optional rotation parameters.

\subsection{Module 2: Intermediate Representation (IR)}
The IR module (\texttt{src/qvm/ir.py}) defines the \texttt{QuantumCircuit} class, which acts as the single unified data structure across all compiler and simulator passes.

\textbf{Structure of \texttt{QuantumCircuit}:}
\begin{itemize}[leftmargin=*]
    \item \texttt{num\_qubits (int):} The total number of declared logical qubits.
    \item \texttt{operations (list[dict]):} An ordered sequence of operation dictionaries. Each dictionary contains:
    \begin{itemize}
        \item \texttt{name:} Gate or directive string (e.g., \texttt{"h"}, \texttt{"cx"}, \texttt{"rz"}, \texttt{"measure"}).
        \item \texttt{qubits:} List of target qubit indices.
        \item \texttt{params:} List of floating-point parameters (such as rotation angles $\theta$).
        \item \texttt{condition:} Optional classical register guard dictionary \texttt{\{"register": name, "index": idx, "value": val\}}.
        \item \texttt{target\_bit:} Optional tuple \texttt{(register\_name, index)} specifying where measurement outcomes are saved.
    \end{itemize}
    \item \texttt{classical\_registers (dict[str, int]):} Maps register names to their bit widths.
\end{itemize}

\textbf{Key Design Decisions in the IR:}
Early in development, we considered representing classical bits as auxiliary quantum states in the Hilbert space. However, adding $m$ classical bits to an $n$-qubit statevector would blow up the state dimension to $2^{n+m}$, which quickly exhausts memory. Instead, we completely decoupled classical memory from the quantum state vector. Classical registers are represented as a separate integer bit-array dictionary. This allows evaluating classical conditionals (\texttt{if (c == 1) X q;}) at near-zero computational cost.

\subsection{Module 3: Transpiler and Routing Subsystem}
The transpiler (\texttt{src/qvm/transpiler.py}) compiles hardware-agnostic circuits to target physical hardware topologies by inserting SWAP operations.

\textbf{Routing Mechanics:}
The transpiler tracks the evolving mapping between logical qubits and physical processor nodes using two dictionaries: \texttt{logical\_to\_physical} and \texttt{physical\_to\_logical}. Single-qubit gates fire immediately on whatever physical node currently holds the target logical qubit. Two-qubit gates (such as CNOT), however, require checking physical graph adjacency.

If physical qubits $p_1$ and $p_2$ share a coupler in \texttt{TargetArchitecture}, the gate is emitted directly. If they do not, the router intervenes using one of two strategies:
\begin{itemize}[leftmargin=*]
    \item \textbf{Greedy BFS Router:} Uses Breadth-First Search on the hardware graph to find the shortest path between $p_1$ and $p_2$. It inserts a sequence of SWAP gates along the path to move the qubits into adjacency, executes the two-qubit gate, and optionally inserts reverse SWAPs if mapping restoration is active.
    \item \textbf{SABRE Lookahead Router \cite{li2019sabre}:} Maintains an active ready queue of two-qubit gates whose dependencies have been satisfied. For every candidate SWAP on the physical coupling graph, SABRE evaluates an exponential decay cost across upcoming gates:
    \begin{equation}
    \text{Cost}(s) = \frac{1}{|F|} \sum_{g \in F} \text{dist}(p_1^g, p_2^g) + W \cdot \frac{1}{|E|} \sum_{g' \in E} \text{dist}(p_1^{g'}, p_2^{g'})
    \end{equation}
    where $F$ is the front layer of gates, $E$ is the extended lookahead window, and $W = 0.5$ is a weighting factor. This lookahead avoids local minima and drops total SWAP counts by 30\% to 40\% compared to greedy routing.
\end{itemize}

Figure \ref{fig:sabre_lookahead} visualizes this SABRE lookahead routing mechanism.

\begin{figure}[!htb]
\centering
\resizebox{0.90\textwidth}{!}{%
\begin{tikzpicture}[
    box/.style={rectangle, draw=blue!80!black, fill=blue!8, rounded corners=3pt, font=\footnotesize, align=center, inner sep=5pt},
    queuebox/.style={rectangle, draw=orange!80!black, fill=orange!10, rounded corners=3pt, font=\footnotesize, align=center, inner sep=5pt},
    hwbox/.style={circle, draw=black!80, fill=gray!15, font=\footnotesize\bfseries, minimum size=0.8cm},
    hwline/.style={draw=black!70, line width=1.5pt},
    flow/.style={-{Stealth[scale=0.95]}, line width=1pt, draw=black!75}
]

% Hardware Graph (Left)
\node[font=\bfseries\small, text=blue!90!black] at (-4.5, 2.2) {Physical Coupling Graph};
\node[hwbox] (q0) at (-5.8, 1.0) {$P_0$};
\node[hwbox] (q1) at (-4.5, 1.0) {$P_1$};
\node[hwbox] (q2) at (-3.2, 1.0) {$P_2$};
\node[hwbox] (q3) at (-4.5, -0.5) {$P_3$};
\node[hwbox] (q4) at (-3.2, -0.5) {$P_4$};

\draw[hwline] (q0) -- (q1);
\draw[hwline] (q1) -- (q2);
\draw[hwline] (q1) -- (q3);
\draw[hwline] (q2) -- (q4);
\draw[hwline] (q3) -- (q4);

% Lookahead Queues (Right)
\node[font=\bfseries\small, text=orange!90!black] at (2.5, 2.2) {SABRE Heuristic Evaluation};

\node[queuebox, text width=7.0cm] (front) at (2.5, 1.1) {
  \textbf{Front Layer $F$ (Active 2-Qubit Gates)}\\[2pt]
  $g_1: CX(q_0, q_2)$ \quad $\to \text{dist}(P_0, P_2) = 2$ \quad (Blocked!)
};

\node[queuebox, text width=7.0cm, fill=orange!5] (extended) at (2.5, -0.4) {
  \textbf{Extended Lookahead Window $E$ (Upcoming Gates)}\\[2pt]
  $g_2: CX(q_0, q_3)$ \quad $g_3: CX(q_2, q_4)$ \quad $g_4: CX(q_1, q_4)$\\[2pt]
  \textit{Evaluated with exponential distance decay weight $W = 0.5$}
};

\node[box, text width=7.0cm, fill=green!10, draw=green!70!black] (decision) at (2.5, -2.0) {
  \textbf{Optimal Candidate SWAP Selection:}\\[2pt]
  Candidate SWAP $(P_0, P_1)$ yields lowest total cost $\text{Cost}(s)$.\\
  $\Longrightarrow$ Insert $\text{SWAP}(P_0, P_1)$, update mapping, and proceed.
};

\draw[flow] (front) -- (extended);
\draw[flow] (extended) -- (decision);

\end{tikzpicture}%
}
\caption{SABRE Lookahead Routing and Heuristic Evaluation}
\label{fig:sabre_lookahead}
\end{figure}

\subsection{Module 4: Dual Simulation Engines}

\textbf{Statevector Simulator (\texttt{src/qvm/simulator.py}):}
The statevector simulator represents an $n$-qubit quantum state as a 1D NumPy complex128 array of length $2^n$. Naive simulation constructs full $2^n \times 2^n$ unitary matrices via Kronecker products, which exhausts memory at modest qubit counts (a 12-qubit matrix requires $16,777,216$ complex numbers, or 268~MB).

To achieve high execution speed, our simulator avoids building $2^n \times 2^n$ matrices. Instead, we reshape the statevector into an $n$-dimensional tensor of shape $(2, 2, \dots, 2)$. To apply a single-qubit gate $U$ to qubit $k$:
\begin{enumerate}[leftmargin=*]
    \item Reshape $|\psi\rangle$ into shape $(2, 2, \dots, 2)$ with $n$ axes.
    \item Use \texttt{numpy.tensordot} or axis swapping to contract the $2 \times 2$ gate matrix $U$ directly along axis $k$.
    \item Permute the axes back to their canonical order and flatten into $2^n$.
\end{enumerate}
This operation executes in microseconds and requires zero full-matrix allocations. For CNOT and SWAP gates, we use bitwise index slicing and boolean masks, swapping amplitude slices directly in memory.

\textbf{Matrix Product State (MPS) Simulator (\texttt{src/qvm/mps\_simulator.py}):}
To simulate circuits with 20+ qubits, we implemented a Matrix Product State simulator \cite{vidal2003mps}. An $n$-qubit state is factored into a 1D chain of rank-3 tensors:
\begin{equation}
|\psi\rangle = \sum_{s_1, \dots, s_n} A^{[1] s_1} A^{[2] s_2} \dots A^{[n] s_n} |s_1 s_2 \dots s_n\rangle
\end{equation}
where $s_i \in \{0, 1\}$ represents the physical qubit basis state, and the internal indices have dimension at most $\chi$ (the bond dimension).

When applying a two-qubit gate across adjacent tensors $A^{[i]}$ and $A^{[i+1]}$:
\begin{enumerate}[leftmargin=*]
    \item Contract $A^{[i]}$ and $A^{[i+1]}$ along their shared bond index into a two-site tensor $\Theta$ of shape $(\chi_L, 2, 2, \chi_R)$.
    \item Apply the $4 \times 4$ two-qubit gate matrix to the physical indices.
    \item Reshape $\Theta$ into a matrix of shape $(2\chi_L, 2\chi_R)$ and perform Singular Value Decomposition: $\Theta = U \Sigma V^\dagger$.
    \item Truncate the singular value spectrum to bond dimension $\chi$, discarding eigenvalues below a threshold $\epsilon = 10^{-12}$.
    \item Reconstruct updated tensors $A'^{[i]}$ and $A'^{[i+1]}$.
\end{enumerate}

Figure \ref{fig:mps_svd} illustrates this two-site MPS tensor contraction and SVD truncation pipeline.

\begin{figure}[!htb]
\centering
\resizebox{0.92\textwidth}{!}{%
\begin{tikzpicture}[
    tensor/.style={rectangle, draw=blue!80!black, fill=blue!12, rounded corners=3pt, font=\small\bfseries, minimum size=0.9cm},
    gate/.style={rectangle, draw=purple!80!black, fill=purple!15, rounded corners=3pt, font=\small\bfseries, minimum width=2.4cm, minimum height=0.7cm},
    leg/.style={draw=black!80, line width=1.4pt},
    arrow/.style={-{Stealth[scale=1.0]}, line width=1.1pt, draw=black!75}
]

% Initial MPS Chain
\node[font=\bfseries\small, text=blue!90!black] at (-4.5, 2.0) {1D Tensor Train Representation};
\node[tensor] (A1) at (-6.5, 0.8) {$A_1$};
\node[tensor] (A2) at (-4.5, 0.8) {$A_2$};
\node[tensor] (A3) at (-2.5, 0.8) {$A_3$};
\node[tensor] (A4) at (-0.5, 0.8) {$A_4$};

\draw[leg] (A1) -- node[above, font=\tiny]{$\chi$} (A2);
\draw[leg] (A2) -- node[above, font=\tiny]{$\chi$} (A3);
\draw[leg] (A3) -- node[above, font=\tiny]{$\chi$} (A4);

\draw[leg] (A1) -- ++(0, -0.9) node[below, font=\tiny]{$s_1$};
\draw[leg] (A2) -- ++(0, -0.9) node[below, font=\tiny]{$s_2$};
\draw[leg] (A3) -- ++(0, -0.9) node[below, font=\tiny]{$s_3$};
\draw[leg] (A4) -- ++(0, -0.9) node[below, font=\tiny]{$s_4$};

% Gate Application & SVD Truncation (Right)
\node[font=\bfseries\small, text=purple!90!black] at (4.5, 2.0) {Two-Site Contraction \& SVD Truncation};

\node[gate] (U) at (4.5, 0.9) {Two-Qubit Gate $U_{23}$};
\node[tensor] (A2p) at (3.5, -0.5) {$A_2$};
\node[tensor] (A3p) at (5.5, -0.5) {$A_3$};

\draw[leg] (A2p) -- node[above, font=\tiny]{bond} (A3p);
\draw[leg] (U.south -| A2p) -- (A2p);
\draw[leg] (U.south -| A3p) -- (A3p);

\node[rectangle, draw=orange!80!black, fill=orange!10, rounded corners=3pt, font=\footnotesize, text width=6.2cm, align=center] (svd) at (4.5, -2.0) {
  \textbf{Singular Value Decomposition (SVD):}\\
  $\Theta = U \cdot \Sigma \cdot V^\dagger \quad \longrightarrow \quad \text{Truncate to bond } \chi$\\
  \textit{Accumulate truncation error: } $\Delta = \sum_{j > \chi} \sigma_j^2$
};

\draw[arrow] (A2p.south) ++(1.0, -0.2) -- (svd.north);

\end{tikzpicture}%
}
\caption{MPS Tensor Train Architecture and SVD Truncation}
\label{fig:mps_svd}
\end{figure}

\subsection{Module 5: Noise Modeling Engine}
Real quantum computers are susceptible to environmental decoherence. In \texttt{src/qvm/noise.py}, we implemented open-system quantum dynamics using the Kraus operator formalism:
\begin{equation}
\rho' = \mathcal{E}(\rho) = \sum_{k} E_k \rho E_k^\dagger, \quad \text{where } \sum_k E_k^\dagger E_k = I
\end{equation}

Figure \ref{fig:kraus_noise} illustrates the noise channels supported in QVM.

\begin{figure}[!htb]
\centering
\resizebox{0.92\textwidth}{!}{%
\begin{tikzpicture}[
    channel/.style={rectangle, draw=blue!80!black, fill=blue!8, rounded corners=4pt, font=\footnotesize, text width=5.8cm, inner sep=6pt},
    title/.style={font=\bfseries\small, text=blue!90!black}
]

\node[channel] (depol) at (-3.5, 1.5) {
  \textbf{Depolarizing Noise Channel}\\[2pt]
  Applies uniform random Pauli errors with probability $p$:\\
  $\mathcal{E}(\rho) = (1-p)\rho + \frac{p}{3}(X\rho X + Y\rho Y + Z\rho Z)$
};

\node[channel] (t1) at (3.5, 1.5) {
  \textbf{Amplitude Damping ($T_1$ Relaxation)}\\[2pt]
  Models energy decay from $|1\rangle$ to $|0\rangle$:\\
  $E_0 = \begin{pmatrix} 1 & 0 \\ 0 & \sqrt{1-\gamma} \end{pmatrix}, \quad E_1 = \begin{pmatrix} 0 & \sqrt{\gamma} \\ 0 & 0 \end{pmatrix}$\\
  where $\gamma = 1 - e^{-t/T_1}$
};

\node[channel] (t2) at (-3.5, -1.5) {
  \textbf{Phase Damping ($T_2$ Dephasing)}\\[2pt]
  Models loss of quantum phase coherence:\\
  $E_0 = \begin{pmatrix} 1 & 0 \\ 0 & \sqrt{1-\lambda} \end{pmatrix}, \quad E_1 = \begin{pmatrix} 0 & 0 \\ 0 & \sqrt{\lambda} \end{pmatrix}$\\
  where $\lambda = 1 - e^{-t/T_\phi}$
};

\node[channel] (readout) at (3.5, -1.5) {
  \textbf{Asymmetric Readout Error}\\[2pt]
  Applies classical confusion matrix upon measurement:\\
  $M = \begin{pmatrix} 1 - p(1|0) & p(0|1) \\ p(1|0) & 1 - p(0|1) \end{pmatrix}$\\
  Perturbs sampled bitstrings according to detector bias.
};

\end{tikzpicture}%
}
\caption{Kraus Operator Noise Channels in QVM}
\label{fig:kraus_noise}
\end{figure}

\subsection{Module 6: Variational Quantum Algorithms (VQA)}
To support hybrid quantum-classical computing, QVM incorporates variational tooling in \texttt{src/qvm/vqe.py}, \texttt{src/qvm/qaoa.py}, and \texttt{src/qvm/gradient.py}:
\begin{itemize}[leftmargin=*]
    \item \textbf{Analytic Parameter-Shift Rule:} Computes exact analytical gradients for parameterized ansätze without numerical finite-difference approximation errors:
    \begin{equation}
    \frac{\partial \langle H \rangle}{\partial \theta_k} = \frac{1}{2} \left[ \langle H \rangle_{\theta_k + \frac{\pi}{2}} - \langle H \rangle_{\theta_k - \frac{\pi}{2}} \right]
    \end{equation}
    \item \textbf{VQE Solver:} Uses classical optimizers (COBYLA, BFGS, Nelder-Mead) to find ground-state eigenvalues for molecular Hamiltonians (verified on molecular $H_2$).
    \item \textbf{QAOA Solver:} Implements the Quantum Approximate Optimization Algorithm for graph combinatorial problems like MaxCut using alternating cost and mixer unitary layers.
\end{itemize}

\section{Algorithm Implementation}

\subsection{SABRE Routing Algorithm}
Below is the core loop of the SABRE lookahead routing heuristic implemented in \texttt{src/qvm/transpiler.py}:

\begin{algorithm}[H]
\caption{SABRE Routing with Decay-Weighted Lookahead}
\label{alg:sabre}
\begin{algorithmic}[1]
\Require Logical circuit $\mathcal{C}$, Target coupling graph $G=(V, E)$, Lookahead weight $W$, Decay factor $\delta$
\Ensure Physical circuit $\mathcal{C}_{phys}$ with valid topological SWAPs
\State Initialize qubit mapping $\pi: Q_{logical} \to V_{physical}$ as identity
\State Compute all-pairs shortest path distance matrix $D$ on $G$
\State Initialize ready queue $F$ with gates whose dependencies are met
\State Initialize extended lookahead queue $E$ with future gates
\While{$F$ is not empty}
    \State Extract gate $g = (q_1, q_2)$ from $F$
    \If{$( \pi(q_1), \pi(q_2) ) \in E$}
        \State Emit $g$ to $\mathcal{C}_{phys}$ mapped to $(\pi(q_1), \pi(q_2))$
        \State Advance dependency queues $F$ and $E$
    \Else
        \State $\text{best\_swap} \leftarrow \text{None}, \quad \text{min\_cost} \leftarrow \infty$
        \For{each candidate physical edge $(s_1, s_2) \in E$}
            \State Temporarily apply swap $(s_1, s_2)$ to $\pi$
            \State $\text{Cost} \leftarrow \frac{1}{|F|}\sum_{f \in F} D(\pi(f_1), \pi(f_2)) + W \sum_{e \in E} \delta^{\text{depth}(e)} D(\pi(e_1), \pi(e_2))$
            \If{$\text{Cost} < \text{min\_cost}$}
                \State $\text{min\_cost} \leftarrow \text{Cost}, \quad \text{best\_swap} \leftarrow (s_1, s_2)$
            \EndIf
            \State Revert temporary swap $(s_1, s_2)$ in $\pi$
        \EndFor
        \State Emit $\text{SWAP}(\text{best\_swap})$ to $\mathcal{C}_{phys}$
        \State Update mapping $\pi$ permanently with $\text{best\_swap}$
    \EndIf
\EndWhile
\State \Return $\mathcal{C}_{phys}$
\end{algorithmic}
\end{algorithm}

\section{External APIs and SDKs}
Table \ref{tab:external_libs} summarizes the external libraries integrated into QVM.

\begin{table}[h!]
\centering
\caption{External Libraries and Integration Points}
\label{tab:external_libs}
\renewcommand{\arraystretch}{1.3}
\begin{tabularx}{\textwidth}{|p{3cm}|p{2cm}|X|}
\hline
\textbf{Library} & \textbf{Version} & \textbf{Role in QVM} \\
\hline
NumPy & 1.24+ & Core linear algebra, complex128 statevector manipulation, tensor slicing \\
\hline
Lark & 1.1+ & Formal EBNF parsing for OpenQASM 3.0 language specifications \\
\hline
FastAPI & 0.100+ & High-performance asynchronous REST API service \\
\hline
Matplotlib & 3.7+ & Generation of publication-quality circuit schematics and probability histograms \\
\hline
pytest & 7.4+ & Test runner for unit, integration, and algorithmic correctness suites \\
\hline
\end{tabularx}
\end{table}

\section{User Interface Implementation}

\subsection{Command-Line Interface (CLI)}
The CLI (\texttt{src/qvm/cli.py}) exposes all compiler and simulator features through a clean command-line interface.

\textbf{Execution Example:}
\begin{verbatim}
# Run a circuit with SABRE transpilation and plotting
python -m src.qvm.cli examples/grover_101.json --nqubits 3 \
    --transpile --routing sabre --shots 1000 --visualize
\end{verbatim}

Figure \ref{fig:cli_execution} displays terminal output from the CLI running canonical quantum algorithm examples, showing execution duration, measurement counts, and classical memory states.

\begin{figure}[!htb]
\centering
\includegraphics[width=0.95\textwidth,height=0.8\textheight,keepaspectratio]{../ScreenShots/CLI_Output.png}
\caption{CLI Execution Output --- Algorithm Verification Traces}
\label{fig:cli_execution}
\end{figure}

\subsection{Web API}
The FastAPI service (\texttt{api/app.py}) provides REST endpoints for automated integration.

\textbf{Primary Endpoints:}
\begin{itemize}[leftmargin=*]
    \item \texttt{POST /run:} Accepts a JSON payload containing circuit code (QASM or JSON), optional transpilation flags, noise parameters, and shot counts. Returns execution status, statevector amplitudes, measurement probabilities, and duration.
    \item \texttt{GET /health:} Returns service health and version metadata.
    \item \texttt{GET /circuits:} Queries saved circuits from Supabase cloud database.
    \item \texttt{GET /history:} Retrieves past execution records and performance metrics.
\end{itemize}

\subsection{Web Dashboard (Next.js 16 \& React 19)}
The web frontend is a single-page application built with Next.js 16, React 19, TypeScript, and Tailwind CSS 4, deployed on Vercel.

Figures \ref{fig:web_gui_interface} and \ref{fig:web_gui_results} illustrate the web user interface in action:
\begin{itemize}[leftmargin=*]
    \item Figure \ref{fig:web_gui_interface} shows the main circuit editing interface with syntax-highlighted editor, preset algorithm selectors, and transpilation controls.
    \item Figure \ref{fig:web_gui_results} displays execution results, including the rendered SVG circuit schematic and interactive measurement probability distribution.
\end{itemize}

\begin{figure}[!htb]
\centering
\includegraphics[width=0.95\textwidth,height=0.8\textheight,keepaspectratio]{../ScreenShots/RuningExapmle_WebUI01.png}
\caption{Web GUI Main Interface --- Circuit Editor and Configuration Panel}
\label{fig:web_gui_interface}
\end{figure}

\begin{figure}[!htb]
\centering
\includegraphics[width=0.95\textwidth,height=0.8\textheight,keepaspectratio]{../ScreenShots/RuningExapmle_WebUI02.png}
\caption{Web GUI Results Display --- Circuit Schematic and Outcome Probabilities}
\label{fig:web_gui_results}
\end{figure}

\section{Deployment and Installation}

\subsection{Local Installation}
QVM installs directly via standard Python tooling:
\begin{verbatim}
# 1. Clone repository
git clone https://github.com/qayum/quantum-virtual-machine.git
cd quantum-virtual-machine

# 2. Create and activate virtual environment
python -m venv .venv
# On Windows: .venv\Scripts\activate
# On Linux/macOS: source .venv/bin/activate

# 3. Install dependencies
pip install -r requirements.txt

# 4. Verify test suite
python -m pytest tests/
\end{verbatim}

\subsection{Running the Local API Server}
\begin{verbatim}
# Launch ASGI web server on port 8000
python -m src.qvm.server --host 127.0.0.1 --port 8000
\end{verbatim}

\section{Chapter Summary}
This chapter detailed the concrete implementation of the Quantum Virtual Machine. We examined the Lark-based OpenQASM 3.0 two-pass parser, the \texttt{QuantumCircuit} intermediate representation, and the tensor-slicing optimizations powering the statevector simulator. We presented the SABRE lookahead routing algorithm and the Matrix Product State (MPS) tensor network engine with SVD bond truncation. We added publication-quality vector diagrams illustrating SABRE heuristic evaluation, MPS tensor contractions, and Kraus noise channels. Finally, we reviewed the CLI and Next.js web dashboard implementations. Chapter 6 presents the testing and evaluation methodology that verified the correctness of these components.
